\subsection{和と差}
微分して元の関数になることを,数列の和でも考えてみよう.\,$S_n=\sum_{k=1}^n a_k$なら,
\[
 S_n-S_{n-1}=a_n
\]
である.和の隣り合う二つを引くと,最後に加えた一項が残る.一方,積分では
\[
 A(x+h)-A(x)=\int_x^{x+h}f(t)\dd t
\]
となり,新しく加えた短い区間が残る.\,$h$で割って極限を取れば,\,$A'(x)=f(x)$である.数列では一項分の差を取り,関数では短い区間の変化率を考えるという違いはあるが,いずれも,足したものからその場所で加わった量を取り出していると読める.

例えば速度$v(t)=2t$なら,原始関数は$t^2$なので,\,$\int_0^T2t\dd t=T^2$である.導入では長方形の和と三角形の面積から求めた値が,今度は微分の公式から得られた.一般の速度でも,原始関数が見つかれば同じ計算ができる.

